
% LaTeX source for expanded academic paper
\documentclass[12pt]{article}
\usepackage{amsmath,amssymb,amsfonts,geometry,setspace}
\geometry{margin=1in}
\onehalfspacing
\title{Relativity Does Not Contain Classical Mechanics:\\Geometry, Time, Constraints, and the Breakdown of Newtonian Structure}
\author{}
\date{}

\begin{document}
\maketitle

\begin{abstract}
It is widely asserted that Newtonian mechanics is recovered from special relativity in the limit of small velocities or formally in the limit $c \to \infty$. This pedagogical belief is conceptually mistaken. Classical mechanics is not contained within relativity, not even as a limiting case. The two theories differ fundamentally in their geometric structure, their treatment of time, their dynamical principles, and the ontology of multi-particle systems. Without a global time function, the canonical Hamiltonian structure of Newtonian mechanics collapses. Instantaneous constraints, rigid bodies, and the Newtonian configuration-space framework cannot be formulated in relativistic spacetime. No-interaction theorems further demonstrate that interacting relativistic particle mechanics is impossible without fields. Only the motion of a single particle under an external electromagnetic field survives intact. This paper develops a unified geometric and dynamical argument demonstrating the categorical failure of the correspondence between classical and relativistic mechanics, clarifying what aspects of Newtonian theory can be approximated and what structures are irrecoverably lost.
\end{abstract}

\section{Introduction}
The conventional wisdom that classical mechanics is recovered from relativity when $v/c \ll 1$ is deeply entrenched in the physics literature. Textbooks frequently state that Newtonian formulas emerge by taking the first terms in a Taylor expansion of relativistic expressions, or by taking the limit $c \to \infty$. While such approximations correctly reproduce certain numerical results, they do not recover the \emph{theory} of classical mechanics. A physical theory is more than formulas; it is a geometric and dynamical structure. The absence of absolute simultaneity, the lack of global time, incompatibility of instantaneous constraints, and the necessity of fields for interactions signal a profound structural divergence.

The aim of this paper is to synthesize several independent lines of analysis into a unified argument: classical mechanics is not a special case of relativity. What relativity recovers is a restricted kinematic approximation, not the Newtonian world with its rigid bodies, constraints, absolute time, and multi-particle Hamiltonian evolution.

\section{Geometric Foundations: Galilean vs. Lorentzian Structure}
Newtonian mechanics is inseparable from Galilean spacetime, which possesses absolute time and Euclidean space. The geometry consists of a three-dimensional spatial manifold evolving in a universal time $t$. Simultaneity is observer-independent. The symmetry group is the Galilei group, whose Lie algebra includes a central extension (mass).

In contrast, special relativity is formulated in Minkowski spacetime with metric signature $(-,+,+,+)$. There is no invariant simultaneity, and the symmetry group is the Poincar'e group. The light-cone structure imposes causal limitations incompatible with instantaneous interactions.

No smooth geometric limit deforms Lorentzian spacetime into Galilean spacetime without singularities. The limit $c \to \infty$ collapses the light cone into a spatial hyperplane, destroying the causal order. Group-theoretically, the passage from Poincar'e to Galilei symmetry is an In"on"u--Wigner contraction, not a geometric limit. This contraction loses essential structural information.

Thus relativity does not contain Galilean spacetime in any continuous manner.

\section{The Role of Time and the Collapse of Hamiltonian Structure}
\subsection{Global time in classical mechanics}
Newtonian mechanics requires a global time parameter shared by all degrees of freedom. This time orders events, defines simultaneity, and underlies the Hamiltonian formulation. Phase space trajectories evolve with respect to a universal time coordinate. The Newtonian $N$-particle configuration space $\mathbb{R}^{3N}$ is defined on simultaneous positions of all particles.

\subsection{Lack of global time in relativity}
Special relativity forbids any invariant foliation of spacetime. Simultaneity is frame-dependent. General relativity goes further: generic spacetimes lack any global time function. Even globally hyperbolic spacetimes admit many inequivalent foliations, none of which is privileged by the theory.

Without global time:
\begin{itemize}
    \item there is no universal evolution parameter for many-particle systems;
    \item canonical Hamiltonian mechanics depends on the arbitrary choice of a foliation;
    \item configuration space does not exist in the Newtonian sense;
    \item instantaneous constraints cannot be formulated.
\end{itemize}

The Newtonian phase-space structure is destroyed at a foundational level.

\section{Constraints and the Impossibility of Rigid Bodies}
Rigid bodies in classical mechanics require instantaneous constraint enforcement and a notion of shape preserved by all inertial observers. These presuppose global simultaneity. In relativity, Born rigidity is possible only under highly restricted motions and cannot be maintained for general accelerations. Constraint forces propagate at finite speed, so multi-point instantaneous constraints are prohibited.

Constraints defining rigid links, rods, joints, or holonomic relations are not relativistically admissible, except as approximations mediated through elastic or electromagnetic fields with finite propagation speed. The configuration spaces associated with classical mechanical systems (e.g., rigid rotors, articulated bodies) have no relativistic analogue.

\section{Multi-Particle Dynamics and the No-Interaction Theorems}
The Currie--Jordan--Sudarshan no-interaction theorem demonstrates that a system of finitely many relativistic particles cannot interact through potentials if one demands Poincar'e invariance and a single time parameter. Nontrivial interactions require fields. Thus classical particle mechanics with interparticle forces has no relativistic extension.

This result explains the structural collapse of the Newtonian $N$-body problem in relativity. The notion of an instantaneous force depending on simultaneous positions is meaningless. Relativity allows only local interactions mediated by fields.

\section{Relativistic Continuum Mechanics}
One might hope that classical continuum mechanics survives relativity even if rigid bodies do not. However, relativistic continuum mechanics is fundamentally different from its Newtonian counterpart. A relativistic continuum is defined by a stress--energy tensor $T^{\mu\nu}$ satisfying local conservation laws. There is no concept of a perfectly rigid body, and the equations of state must respect causality.

In Newtonian continuum mechanics, pressure and stress can adjust instantaneously across the medium. In relativity, signals propagate at finite speed, invalidating classical constitutive relations. The Newtonian limit of relativistic continua yields compressible, not rigid, media. The structural assumptions are thus incompatible.

\section{What Survives: Single-Particle Motion in External Electromagnetic Fields}
Given the collapse of constraints, multi-particle structure, and continuum rigidity, the only part of classical mechanics that survives intact in relativity is the motion of a single particle under an external electromagnetic field. The Lorentz force law is compatible with Minkowski geometry and requires no global time. Maxwell's equations provide the necessary field mediation.

Everything else---interacting particles, rigid bodies, constrained systems, Newtonian continua---is not contained within relativity.

\section{Conclusion}
Classical mechanics is not a limit case of relativity. The two theories rest on fundamentally different spacetime geometries, temporal structures, and dynamical frameworks. Relativity preserves only certain approximations of Newtonian formulas, not its theoretical structure. The absence of global time, impossibility of instantaneous constraints, collapse of configuration space, and no-interaction theorems demonstrate the incompatibility.

Relativity does not generalize classical mechanics; it replaces it.

\section*{References}
\begin{itemize}
    \item Ehlers, J. ``The Newtonian Limit of General Relativity.'' In: \emph{Classical Mechanics and Relativity}. Springer.
    \item Trautman, A. ``Foundations and Current Problems of General Relativity.'' Brandeis Summer School, 1964.
    \item Currie, D. G., Jordan, T. F., \& Sudarshan, E. C. G. ``Relativistic Invariance and Hamiltonian Theories of Interacting Particles.'' \emph{Rev. Mod. Phys.} 35, 350 (1963).
    \item Born, M. ``Die Theorie des starren Elektrons in der Kinematik des Relativit"atsprinzips.'' \emph{Ann. Phys.} 335, 1 (1909).
    \item In"on"u, E., \& Wigner, E. P. ``On the Contraction of Groups and their Representations.'' \emph{Proc. Natl. Acad. Sci.} 39 (1953).
\end{itemize}

% Supplemental LaTeX: Mathematical derivations and proofs
\section*{Supplement: Mathematical Derivations}
\subsection*{1. Lorentz transformation expansion}
The standard Lorentz transformation (boost in the $x$-direction) is
\begin{align}
 t' &= \gamma\left(t - \frac{v x}{c^2}\right), \\
 x' &= \gamma(x - v t), \\
 y' &= y, \quad z' = z,
\end{align}
with $\gamma=(1-v^2/c^2)^{-1/2}$. For $\varepsilon = v/c \ll 1$ expand
\begin{equation}
 \gamma = 1 + \tfrac12\varepsilon^2 + \tfrac{3}{8}\varepsilon^4 + O(\varepsilon^6).
\end{equation}
Substituting gives the low-velocity expansion
\begin{align}
 t' &= t - \frac{v x}{c^2} + \frac{1}{2}\frac{v^2}{c^2} t + O(\varepsilon^4), \\
 x' &= x - v t + \frac{1}{2}\frac{v^2}{c^2} x + O(\varepsilon^4).
\end{align}
Thus the Galilean transformation $t'=t,\ x'=x-vt$ is recovered only to leading order; higher-order terms modify simultaneity by amounts $\sim v x/c^2$.

\subsection*{2. Metric degeneration and Newton--Cartan fields}
Start from Minkowski metric $g_{\mu\nu}$ with components
\begin{equation}
 g_{00} = -c^2, \quad g_{0i}=0, \quad g_{ij}=\delta_{ij}.
\end{equation}
Define the temporal one-form $\tau_\mu = c^{-1} g_{\mu 0}$ so that $\tau_0 = -1,\ \tau_i=0$. Introduce the spatial contravariant metric
\begin{equation}
 h^{\mu\nu} = \lim_{c\to\infty} g^{\mu\nu} + O(c^{-2}),
\end{equation}
which yields $h^{00}=0,\ h^{0i}=0,\ h^{ij}=\delta^{ij}$. One obtains
\begin{equation}
 h^{\mu\nu}\tau_\nu = 0,
\end{equation}
so $(\tau_\mu,h^{\mu\nu})$ define Newton--Cartan structure. The limiting process is singular because $\det g_{\mu\nu}\to0$ and $g^{\mu\nu}$ ceases to be well-defined uniformly as $c\to\infty$.

\subsection*{3. Algebraic contraction: Poincar\'e to Galilei}
Write the Poincar\'e commutator
\begin{equation}
 [K_i,K_j] = -i\frac{1}{c^2} \epsilon_{ijk} J_k.
\end{equation}
Rescale boosts $\tilde K_i = c K_i$. Then
\begin{equation}
 [\tilde K_i,\tilde K_j] = -i c^2 \frac{1}{c^2} \epsilon_{ijk} J_k = -i \epsilon_{ijk} J_k.
\end{equation}
If instead one takes the limit with unrescaled $K_i$ the right hand side vanishes, producing the Galilei algebra where boosts commute. The mass central extension enters because the projective representations produce an extra term in $[P_i, K_j]$ proportional to $i M \delta_{ij}$.

\subsection*{4. No-interaction theorem (sketch)}
Assume canonical Poisson brackets $\{x^i_a,p_{j,b}\} = \delta^i_j \delta_{ab}$ and attempt to construct Poincar\'e generators as phase-space functions. The requirement that the Poisson brackets of these generators reproduce the Poincar\'e algebra yields functional equations for interaction potentials. One can show these imply vanishing interaction derivatives, forcing potentials to be constant. The detailed algebraic manipulations are in Currie et al. (1963). The conceptual takeaway is that single-time Hamiltonian particle interactions are incompatible with Lorentz invariance.

\subsection*{5. Born rigidity equations}
Using the decomposition of the covariant derivative into kinematic quantities
\begin{equation}
 \nabla_\mu u_\nu = \omega_{\mu\nu} + \sigma_{\mu\nu} + \tfrac{1}{3} \theta h_{\mu\nu} - a_\nu u_\mu,
\end{equation}
with vorticity $\omega_{\mu\nu}$, shear $\sigma_{\mu\nu}$, expansion $\theta$, and four-acceleration $a_\mu$, the Born condition $\mathcal{L}_u h_{\mu\nu}=0$ is equivalent to
\begin{equation}
 2\sigma_{\mu\nu} + \tfrac{2}{3}\theta h_{\mu\nu} = 0, \qquad \omega_{\mu\nu} = 0.
\end{equation}
For generic motion $\sigma_{\mu\nu}$ and $\theta$ are nonzero, so exact Born rigidity fails.

\subsection*{6. ADM constraints and foliation dependence}
The Hamiltonian constraint can be written explicitly
\begin{equation}
 \mathcal{H} = \frac{1}{\sqrt{h}}\left(\pi^{ij}\pi_{ij} - \tfrac{1}{2}(\pi^i{}_i)^2\right) - \sqrt{h} R^{(3)} + 16\pi G \mathcal{H}_{\text{matter}} = 0.
\end{equation}
The momentum constraint enforces spatial diffeomorphism invariance. Solutions require choosing lapse and shift; different choices correspond to distinct slicings related by spacetime diffeomorphisms. No canonical slicing singled out by the dynamics corresponds to Newtonian absolute time.

\subsection*{7. Hyperbolicity and causality}
Linear perturbations of a perfect fluid about a homogeneous background satisfy a wave equation of the form
\begin{equation}
 \Box\Phi + c_s^2 \Delta \Phi = 0,
\end{equation}
with finite sound speed $c_s\le c$, guaranteeing causal propagation. In Newtonian models the elliptic Poisson equation for instantaneous potentials is incompatible with finite-speed causality.

\subsection*{8. Quantum field limit details}
Applying the Foldy--Wouthuysen transformation to the Dirac Hamiltonian yields the nonrelativistic Hamiltonian
\begin{equation}
 H_{\text{NR}} = m + \frac{\mathbf{p}^2}{2m} - \frac{\mathbf{p}^4}{8 m^3} + \cdots,
\end{equation}
but this expansion is performed after selecting a positive-frequency subspace; particle number conservation, absence of pair creation, and a fixed-time description are all extra assumptions that break down once relativistic effects are significant.

\section*{Remarks}
These derivations make explicit the various senses in which Newtonian structures arise only asymptotically or via singular limiting procedures. They underline the paper's central claim: approximation of formulas is not the same as theoretical containment.
\end{document}
